\section{Background and related work}\label{sec:related}

\paragraph{Learning-based vulnerability detection and its evaluation}
Function-level detectors based on graph networks~\cite{zhou2019devign}, token models and pre-trained code language models are commonly reported with high accuracy on random splits of public corpora. Chakraborty et al.~\cite{chakraborty2021deep} showed that performance drops sharply on realistic data, and recent studies with PrimeVul~\cite{ding2024primevul} reach similar conclusions for code language models under de-duplicated, time-aware evaluation. The data underneath is itself a threat: label noise and duplication in vulnerability corpora are documented~\cite{croft2023data}, and code duplication inflates the results of machine-learning models of code~\cite{allamanis2019adverse}. The same concern about space--time bias is well established in malware detection~\cite{pendlebury2019tesseract}. Our audit (Section~\ref{sec:design}) and the leakage experiment (Section~\ref{sec:rq1}) quantify these issues for the two corpora we use and carry the analysis one step further, from how \emph{accurate} a detector is to what a team can \emph{safely delegate} to it.

\paragraph{Calibration, selective prediction and uncertainty}
Neural networks are often miscalibrated~\cite{guo2017calibration} and post-hoc scaling and isotonic regression are standard remedies. Selective classification lets a model abstain on uncertain inputs and trades coverage for risk~\cite{geifman2017selective}. Uncertainty estimation has begun to appear in vulnerability detection, for example Bayesian recurrent detectors and toolkits that quantify the uncertainty of code language models. These approaches provide \emph{scores}, not statistical guarantees on the error of the delegated decisions, and calibration alone does not give a guarantee on a specified class-conditional error such as the share of missed vulnerabilities.

\paragraph{Conformal prediction and risk control}
Conformal prediction~\cite{vovk2005algorithmic,shafer2008tutorial,angelopoulos2021gentle} converts any scorer into a procedure with finite-sample, distribution-free coverage under exchangeability; conformal risk control~\cite{angelopoulos2023conformal} and Learn-then-Test~\cite{angelopoulos2021ltt} generalise it to controlling expected loss and to high-probability statements about arbitrary risks; weighted conformal prediction~\cite{tibshirani2019conformal} extends validity to covariate shift with a known likelihood ratio. Class-conditional miss-rate control of this kind is used for safety warning systems in robotics~\cite{luo2022sample}. To the best of our knowledge, and based on our literature search, the application of such guarantees to \emph{automated vulnerability triage}, together with an empirical analysis of how the guarantee behaves under the project, temporal and cross-dataset shifts that characterise software evolution, has not been studied before. Our contribution is not a new conformal theory, but (i) its specialisation to the clearance decision of a vulnerability detector, with a diagnostic bound on the effect of shift (Proposition~\ref{prop:shift}), the high-probability variant, and an economic rule for choosing the budget; and (ii) a large empirical test of when these guarantees hold in practice.

\paragraph{Static analysis and alert fatigue}
Rule-based analysers produce many alerts that developers ignore~\cite{johnson2013why}. We include Flawfinder as a representative lightweight analyser and compare it with the learned triage layer on identical functions.

\paragraph{Robustness and explanation of code models}
Models of code are sensitive to semantics-preserving transformations such as identifier renaming~\cite{rabin2021generalizability,yefet2020adversarial}. For trustworthy automation what matters is whether the \emph{delegated decision} changes, so we measure decision flips of the triage layer rather than only score changes. For explanation we use TreeSHAP~\cite{lundberg2017shap,lundberg2020treeshap} on an interpretable metric-based model, evaluate fidelity by deletion, and test consistency of the explanations across seeds and across datasets.
